% GENERATED FILE. Edit canonical lessons, then run npm run generate:books.
% canonical-source-hash: fnv1a64:5acd0522f3daf6dc
% canonical-lessons: LA-C278-volumen, LA-C278-pagina, LA-C278-littera, LA-C278-verbum, LA-C278-nomen

\chapter{Book-roll, Page, Letter of the alphabet, Word, Name}
\label{ch:la-book-roll-page-letter-of-the-alphabet-word-name}
\hlchaptermodality{278}

\begin{chapteropening}
\textbf{By the end of this chapter:} \emph{I can say a book-roll, a page, a letter of the alphabet, a word and a name.}

Complete the last lesson of chapter 278: I can say a book-roll, a page, a letter of the alphabet, a word and a name.
\end{chapteropening}

\section[volūmen, volūminis]{volūmen, volūminis — a book-roll}

\label{lesson:LA-C278-volumen}



\begin{quote}
Before the new one: say the Latin for an arrow, then the Latin for a trumpet.
\end{quote}

\subsection*{You'll want to know: volūmen, volūminis}
\textbf{volūmen, volūminis} — \textquotedblleft{}a book-roll\textquotedblright{}.

\begin{cousinweb}
\textbf{volūmen} is built on \emph{volvō}, \textquotedblleft{}to roll\textquotedblright{} (Lewis \& Short), with \emph{-men}, an ending for a result: something rolled up. \emph{volvō} is usually traced to Proto-Indo-European \emph{*welw-}, \textquotedblleft{}to roll\textquotedblright{} (EtymDB 2.1). English \emph{volume} comes from it.
\end{cousinweb}

\begin{grammarlens}[title={What you've built}]
One more word for this chapter.
\end{grammarlens}

\subsection*{Your turn}
\begin{itemize}
\raggedright
  \item \emph{Say it:} \emph{volūmen, volūminis}
  \item \emph{Say it:} \emph{volūmen, volūminis}, once more
  \item \emph{Say it:} say \emph{volūmen, volūminis}
  \item \emph{Recall:} say the Latin for an arrow, then the Latin for a trumpet, then say \emph{volūmen, volūminis} again
\end{itemize}

\begin{tcolorbox}[breakable,colback=teal!4,colframe=teal!35!black,title={Before you move on}]
What does volūmen, volūminis mean? (\textquotedblleft{}A book-roll\textquotedblright{}.) Say it once more.
\end{tcolorbox}
\section[pāgina, pāginae]{pāgina, pāginae — a page}

\label{lesson:LA-C278-pagina}



\begin{quote}
Before the new one: say the Latin for a trumpet, then the Latin for a book-roll.
\end{quote}

\subsection*{You'll want to know: pāgina, pāginae}
\textbf{pāgina, pāginae} — \textquotedblleft{}a page\textquotedblright{}.

\begin{cousinweb}
\textbf{pāgina} is built on the root of \emph{pangō}, \textquotedblleft{}to fasten, to fix\textquotedblright{} (Lewis \& Short): first a row of vines fastened together, then a column of writing. The root is usually traced to Proto-Indo-European \emph{*peh2g-}, \textquotedblleft{}to fasten\textquotedblright{} (EtymDB 2.1).
\end{cousinweb}

\begin{grammarlens}[title={What you've built}]
One more word for this chapter.
\end{grammarlens}

\subsection*{Your turn}
\begin{itemize}
\raggedright
  \item \emph{Say it:} \emph{pāgina, pāginae}
  \item \emph{Say it:} \emph{pāgina, pāginae}, once more
  \item \emph{Say it:} say \emph{pāgina, pāginae}
  \item \emph{Recall:} say the Latin for a trumpet, then the Latin for a book-roll, then say \emph{pāgina, pāginae} again
\end{itemize}

\begin{tcolorbox}[breakable,colback=teal!4,colframe=teal!35!black,title={Before you move on}]
What does pāgina, pāginae mean? (\textquotedblleft{}A page\textquotedblright{}.) Say it once more.
\end{tcolorbox}
\section[littera, litterae]{littera, litterae — a letter of the alphabet}

\label{lesson:LA-C278-littera}



\begin{quote}
Before the new one: say the Latin for a book-roll, then the Latin for a page.
\end{quote}

\subsection*{You'll want to know: littera, litterae}
\textbf{littera, litterae} — \textquotedblleft{}a letter of the alphabet\textquotedblright{}.

\begin{cousinweb}
\textbf{littera} has no sure origin; the old link with \emph{linō}, \textquotedblleft{}to smear\textquotedblright{}, is doubtful, and a loan from Greek \emph{diphthera}, \textquotedblleft{}a hide for writing on\textquotedblright{}, through Etruscan has also been suggested.
\end{cousinweb}

\begin{grammarlens}[title={What you've built}]
One more word for this chapter.
\end{grammarlens}

\subsection*{Your turn}
\begin{itemize}
\raggedright
  \item \emph{Say it:} \emph{littera, litterae}
  \item \emph{Say it:} \emph{littera, litterae}, once more
  \item \emph{Say it:} say \emph{littera, litterae}
  \item \emph{Recall:} say the Latin for a book-roll, then the Latin for a page, then say \emph{littera, litterae} again
\end{itemize}

\begin{tcolorbox}[breakable,colback=teal!4,colframe=teal!35!black,title={Before you move on}]
What does littera, litterae mean? (\textquotedblleft{}A letter of the alphabet\textquotedblright{}.) Say it once more.
\end{tcolorbox}
\section[verbum, verbī]{verbum, verbī — a word}

\label{lesson:LA-C278-verbum}



\begin{quote}
Before the new one: say the Latin for a page, then the Latin for a letter of the alphabet.
\end{quote}

\subsection*{You'll want to know: verbum, verbī}
\textbf{verbum, verbī} — \textquotedblleft{}a word\textquotedblright{}.

\begin{cousinweb}
\textbf{verbum} is inherited, usually traced to Proto-Indo-European \emph{*werdhom}, \textquotedblleft{}word\textquotedblright{} (EtymDB 2.1); English \emph{word} is a cousin (Lewis \& Short).
\end{cousinweb}

\begin{grammarlens}[title={What you've built}]
One more word for this chapter.
\end{grammarlens}

\subsection*{Your turn}
\begin{itemize}
\raggedright
  \item \emph{Say it:} \emph{verbum, verbī}
  \item \emph{Say it:} \emph{verbum, verbī}, once more
  \item \emph{Say it:} say \emph{verbum, verbī}
  \item \emph{Recall:} say the Latin for a page, then the Latin for a letter of the alphabet, then say \emph{verbum, verbī} again
\end{itemize}

\begin{tcolorbox}[breakable,colback=teal!4,colframe=teal!35!black,title={Before you move on}]
What does verbum, verbī mean? (\textquotedblleft{}A word\textquotedblright{}.) Say it once more.
\end{tcolorbox}
\section[nōmen, nōminis]{nōmen, nōminis — a name}

\label{lesson:LA-C278-nomen}



\begin{quote}
Before the new one: say the Latin for a letter of the alphabet, then the Latin for a word.
\end{quote}

\subsection*{You'll want to know: nōmen, nōminis}
\textbf{nōmen, nōminis} — \textquotedblleft{}a name\textquotedblright{}.

\begin{cousinweb}
\textbf{nōmen} is inherited, usually traced to Proto-Indo-European \emph{*h1nomn}, \textquotedblleft{}name\textquotedblright{} (EtymDB 2.1); English \emph{name} is a cousin.
\end{cousinweb}

\begin{grammarlens}[title={What you've built}]
One more word for this chapter.
\end{grammarlens}

\subsection*{Your turn}
\begin{itemize}
\raggedright
  \item \emph{Say it:} \emph{nōmen, nōminis}
  \item \emph{Say it:} \emph{nōmen, nōminis}, once more
  \item \emph{Say it:} say \emph{nōmen, nōminis}
  \item \emph{Recall:} say the Latin for a letter of the alphabet, then the Latin for a word, then say \emph{nōmen, nōminis} again
\end{itemize}

\begin{tcolorbox}[breakable,colback=teal!4,colframe=teal!35!black,title={Before you move on}]
What does nōmen, nōminis mean? (\textquotedblleft{}A name\textquotedblright{}.) Say it once more.
\end{tcolorbox}
